\documentclass[11pt,a4paper]{article}
\usepackage[T1]{fontenc}
\usepackage[utf8]{inputenc}
\usepackage[main=italian,provide=*]{babel}
\usepackage{amsmath,amssymb}
\usepackage{geometry}
\geometry{margin=2.3cm,top=2.5cm,bottom=2.7cm}
\usepackage{booktabs}
\usepackage[table]{xcolor}
\usepackage{tcolorbox}
\tcbuselibrary{skins,breakable}
\usepackage{titlesec}
\usepackage{enumitem}
\usepackage{seqsplit}
\usepackage{needspace}
\usepackage{fancyhdr}
\usepackage{hyperref}
\hypersetup{colorlinks=true,linkcolor=Sepia,citecolor=Sepia,urlcolor=Sepia}

\definecolor{inkblue}{HTML}{1B3A5C}
\definecolor{lightblue}{HTML}{EAF1F8}
\definecolor{accent}{HTML}{B0562A}
\definecolor{softgray}{HTML}{6B6B6B}
\definecolor{okgreen}{HTML}{2E6B3E}
\definecolor{okgreenbg}{HTML}{EAF5EC}
\definecolor{warnamber}{HTML}{9C6B12}
\definecolor{warnbg}{HTML}{FBF1E0}
\definecolor{advisorpurple}{HTML}{5B3A7A}
\definecolor{advisorbg}{HTML}{F1ECF6}
\definecolor{notebar}{HTML}{9C6B12}

\titleformat{\section}
  {\normalfont\Large\bfseries\color{inkblue}}
  {\thesection}{0.6em}{}[{\color{inkblue!40}\titlerule}]
\titlespacing{\section}{0pt}{0.8em}{0.4em}
\titleformat{\subsection}
  {\normalfont\large\bfseries\color{accent}}
  {}{0em}{}
\titlespacing{\subsection}{0pt}{0.5em}{0.2em}

\pagestyle{fancy}
\fancyhf{}
\renewcommand{\headrulewidth}{0.4pt}
\fancyhead[R]{\small\color{softgray}\thepage}
\fancyfoot[C]{\small\color{softgray}Tesi triennale, Samuele Schiavi --- Relatore: Prof.\ Alessandro Chiesa}

\newtcolorbox{keyeq}[1][]{
  colback=lightblue, colframe=inkblue!70, boxrule=0.6pt,
  arc=2pt, left=8pt, right=8pt, top=4pt, bottom=4pt, #1
}
\newtcolorbox{panelbox}[1]{
  colback=white, colframe=softgray!50, boxrule=0.5pt,
  arc=2pt, left=7pt, right=7pt, top=3pt, bottom=3pt,
  fonttitle=\bfseries\color{inkblue}, title={#1}
}
\newtcolorbox{okbox}[1]{
  colback=okgreenbg, colframe=okgreen!60, boxrule=0.6pt,
  arc=2pt, left=7pt, right=7pt, top=3pt, bottom=3pt,
  fonttitle=\bfseries\color{okgreen}, title={#1}
}
\newtcolorbox{warnbox}[1]{
  colback=warnbg, colframe=warnamber!70, boxrule=0.6pt,
  arc=2pt, left=7pt, right=7pt, top=3pt, bottom=3pt,
  fonttitle=\bfseries\color{warnamber}, title={#1}
}
\newtcolorbox{advisorbox}[1]{
  colback=advisorbg, colframe=advisorpurple!60, boxrule=0.6pt,
  arc=2pt, left=7pt, right=7pt, top=4pt, bottom=4pt,
  fonttitle=\bfseries\color{advisorpurple}, title={#1}
}
\newtcolorbox{editnote}{
  colback=white, colframe=white, boxrule=0pt,
  borderline west={2.2pt}{0pt}{notebar},
  arc=0pt, left=10pt, right=4pt, top=2pt, bottom=2pt,
  fontupper=\itshape\small\color{softgray!90!black}
}
\fancyhead[L]{\small\color{softgray}Circuito compatto per il passo di Trotter --- trimero anello}

\title{\vspace{-1.5em}\color{inkblue}Circuito compatto per il Trotter del trimero anello\\[0.1em]
\Large\normalfont\color{softgray}Verifica empirica, e perché la stessa conclusione del dimero vale ancora di più}
\author{Note di lavoro --- Tesi triennale, Samuele Schiavi\\ Relatore: Prof.\ Alessandro Chiesa}
\date{}

\begin{document}
\sloppy
\maketitle
\thispagestyle{fancy}
\vspace{-2em}

Seguito naturale di \texttt{circuito\_compatto\_teoria\_e\_limiti.pdf} (fatto
per il dimero), applicato ora al trimero anello. La domanda è la stessa: il
circuito Trotter esplicito (gate fisicamente etichettati per bond) si può
comprimere? Questo file raccoglie il test fatto, i numeri, e --- come già
per il dimero --- perché il risultato \textbf{non va usato} come sostituto
del circuito esplicito quando si arriverà al rumore (Parte 2).

\section{Il problema, e una differenza di partenza rispetto al dimero}

Il circuito Trotter del trimero (\texttt{trotter\_trimero\_anello.py}) usa,
per un singolo passo esterno, \textbf{15 gate a due qubit nativi}
($3\times R_{XX}+3\times R_{YY}+9\times R_{ZZ}$: tre bond per lo scambio
$H_{ex}$, tre bond per il DM $H_{DM}$, ciascuno scomposto nei rispettivi
blocchi), più 12 Hadamard e 3 $R_z$ per il campo.

\begin{warnbox}{Differenza cruciale rispetto al dimero}
Per il dimero (2 qubit) esiste un teorema chiuso: ogni unitaria in $SU(4)$
si realizza con \textbf{al più 3 CNOT} (decomposizione di Cartan/camera di
Weyl, Vatan--Williams 2004, Vidal--Dawson 2004 --- vedi
\texttt{circuito\_compatto\_teoria\_e\_limiti.pdf} per la bibliografia
completa e verificata). Per il trimero (3 qubit, $SU(8)$) \textbf{non
esiste un risultato altrettanto chiuso}: non è stata reperita né verificata
alcuna formula analoga con un numero minimo universale di CNOT per
un'unitaria a 3 qubit generica. Questo documento riporta quindi una
verifica \textbf{empirica} (cosa produce il transpiler di Qiskit), non un
teorema, e la differenza va dichiarata esplicitamente: non è un dettaglio,
cambia lo status epistemico di ogni numero qui sotto.
\end{warnbox}

\section{Metodo}

A differenza del dimero, per il trimero non è disponibile
\texttt{TwoQubitWeylDecomposition} (specifico di 2 qubit): manca quindi il
"metodo 1" indipendente (coordinate di Weyl) usato là come cross-check.
Resta solo la sintesi via transpiler:

\begin{enumerate}[leftmargin=1.4em,itemsep=0.25em]
\item costruzione del circuito Trotter esplicito (fisico, gate
etichettati per bond) con \texttt{trotter\_trimero\_anello.py};
\item estrazione della sua unitaria $8\times8$ esatta,
\texttt{Operator(qc).data} (Qiskit, non un calcolo indipendente --- ma il
circuito è già stato validato a precisione macchina contro la costruzione
a matrice nella sessione precedente, quindi l'operatore è affidabile);
\item la matrice passata come blocco numerico opaco,
\texttt{qc.unitary(U, [0,1,2])} --- il transpiler non vede più scambio o
DM, solo 64 numeri complessi;
\item sintesi: \texttt{transpile(qc, basis\_gates=['cx','rz','sx','x'],
optimization\_level=3)};
\item verifica di fedeltà (\texttt{process\_fidelity}) fra compresso ed
esplicito, per assicurarsi che la compressione riproduca esattamente la
stessa unitaria e non un'approssimazione.
\end{enumerate}

\section{Risultati}

\subsection*{Un solo passo}

\begin{center}
\renewcommand{\arraystretch}{1.3}
\begin{tabular}{lcc}
\toprule
& gate a due qubit & fedeltà vs esplicito \\
\midrule
circuito fisico (esplicito) & 15 ($R_{XX}{+}R_{YY}{+}R_{ZZ}$ nativi) & --- \\
compresso (\texttt{optimization\_level=3}) & \textbf{18 CNOT} & $1.0000000000$ \\
\bottomrule
\end{tabular}
\end{center}

Fedeltà a precisione macchina: la compressione riproduce esattamente la
stessa unitaria, solo con una scomposizione in gate diversa. Notare che 18
non è necessariamente il minimo raggiungibile (l'ottimizzatore di Qiskit
non garantisce ottimalità globale) --- è solo quello che
\texttt{optimization\_level=3} ha trovato.

\subsection*{L'intero prodotto a $N$ passi --- il punto concettualmente importante}

Come già osservato per il dimero, comprimere l'\textbf{intero} prodotto
$U_\text{Trotter}(t,N)$ (non passo per passo) resta comunque un'unica
unitaria a 3 qubit, quindi si comprime anch'esso a un numero di gate
\textbf{indipendente da $N$}:

\begin{center}
\renewcommand{\arraystretch}{1.3}
\begin{tabular}{ccccc}
\toprule
\rowcolor{inkblue!10}
$N$ & gate 2q nativi (esplicito) & CNOT (compresso) & fedeltà \\
\midrule
1   & 15   & 18 & $1.0$ \\
2   & 30   & 19 & $1.0$ \\
5   & 75   & 19 & $1.0$ \\
20  & 300  & 19 & $1.0$ \\
100 & 1500 & 19 & $1.0$ \\
\bottomrule
\end{tabular}
\end{center}

Il conteggio compresso si stabilizza a $19$ CNOT già da $N=2$ in poi,
verificato fino a $N=100$: identico fenomeno del dimero (lì il valore fisso
era 3), qui più marcato in termini assoluti di risparmio percentuale, dato
che il circuito esplicito cresce linearmente ($15N$ gate) mentre il
compresso resta piatto.

\begin{editnote}
Punto di calibrazione importante: nell'analisi di convergenza (punto di
lavoro $R_0$, $b{=}0.05$, $D{=}1.93$) servono $N\simeq300$--$1000$ per
un'infedeltà sotto l'1\%--0.01\%. A quei valori di $N$ il circuito
esplicito avrebbe $4500$--$15000$ gate a due qubit, contro i $\sim19$ del
compresso: un risparmio ancora più vistoso di quello già visto nel
dimero, e quindi un invito ancora più forte a non usarlo per la ragione
sbagliata (vedi sezione successiva).
\end{editnote}

\section{Perché non si usa per il rumore (stessa conclusione del dimero, rafforzata)}

\begin{warnbox}{Il compromesso $N$-vs-rumore sparirebbe, e di più}
L'idea della Parte 2 è studiare come l'errore di Trotter ($\sim1/N$,
vuole $N$ grande) e l'errore fisico dei gate sotto rumore (scala col
numero di gate, vuole $N$ piccolo) si bilancino. Usando il circuito
compresso, il costo in gate resta fisso a $\sim19$ \textbf{qualunque sia
$N$}: il compromesso sparisce dal grafico. Con gli $N$ dell'ordine delle
centinaia/migliaia che servono qui (più alti che nel dimero, per via dei
due livelli di bond-splitting e di $D/J$ non piccolo), l'effetto sarebbe
ancora più ingannevole: si misurerebbe il rumore su un'unitaria a
\emph{costo fisico completamente slegato} dal numero di passi realmente
usati per la fisica.
\end{warnbox}

Vale inoltre, identica, l'osservazione già fatta per il dimero: ottenere
il circuito compresso richiede il calcolo classico preliminare
dell'intera matrice $8\times8$ (o della sua potenza $N$-esima) --- fattibile
qui perché il trimero resta piccolo, ma concettualmente opposto allo
spirito di Trotter (evitare la diagonalizzazione/esponenziazione
classica di matrici che diventano intrattabili per sistemi grandi).

\subsection*{Dove la forma compatta avrebbe senso}

Come per il dimero: se in futuro serve rispondere a "qual è il costo
minimo in gate per realizzare \emph{questa specifica} evoluzione a un $t$
fissato" (domanda a sé, legittima, per esempio come riferimento per un
run su hardware reale), il circuito compresso è lo strumento giusto ---
ma va tenuto separato dal test di Trotter-vs-rumore, che deve usare il
circuito esplicito o comunque uno il cui conteggio di gate scali con $N$.

\section{Riepilogo}

\begin{enumerate}[leftmargin=1.4em,itemsep=0.3em]
\item Per il trimero (3 qubit) non esiste un teorema di compressione
universale altrettanto chiuso di quello a 2 qubit del dimero: la verifica
qui è empirica (transpiler), non un minimo dimostrato.
\item Un singolo passo si comprime da 15 gate nativi a 18 CNOT, fedeltà
$1.0$.
\item L'intero prodotto a $N$ passi si comprime a $\sim19$ CNOT
\textbf{indipendentemente da $N$} (verificato fino a $N=100$) --- stesso
fenomeno del dimero, più marcato in valore assoluto.
\item Conclusione invariata: non usare il circuito compresso per lo
studio Trotter-vs-rumore della Parte 2 (farebbe sparire la scalata in $N$
del costo in gate, il fenomeno stesso da studiare); legittimo solo per la
domanda separata del costo minimo a $t$ fissato.
\end{enumerate}

\needspace{15\baselineskip}
\section*{Appendice: comandi Qiskit usati per la verifica}
\addcontentsline{toc}{section}{Appendice}

\begin{panelbox}{Codice}
\small
\begin{verbatim}
from qiskit import QuantumCircuit, transpile
from qiskit.quantum_info import Operator, process_fidelity
from trotter_trimero_anello import trotter_circuit

qc_step = trotter_circuit(J, Jp, b, D, tau, 1)      # un passo esplicito
U = Operator(qc_step).data                          # unitaria 8x8 esatta

qc_compact = QuantumCircuit(3)
qc_compact.unitary(U, [0, 1, 2])                    # blocco opaco
qc_t = transpile(qc_compact, basis_gates=['cx','rz','sx','x'],
                  optimization_level=3)

print(qc_t.count_ops())
print(process_fidelity(Operator(qc_t), Operator(qc_step)))

# stesso procedimento ripetuto su trotter_circuit(J, Jp, b, D, T, N)
# per N = 1, 2, 5, 20, 100, per vedere se il conteggio compresso
# resta fisso al crescere di N (vedi tabella sopra)
\end{verbatim}
\end{panelbox}

\end{document}
